\documentclass[11pt,a4paper]{article}
\usepackage[margin=1in]{geometry}
\usepackage{amsmath,amssymb}
\usepackage{graphicx}
\graphicspath{{../figures/}}
\usepackage{booktabs}
\usepackage{array}
\usepackage{tabularx}
\usepackage{xcolor}
\usepackage[hidelinks]{hyperref}
\usepackage{caption}
\captionsetup{font=small,labelfont=bf}

\newcommand{\bp}{\ensuremath{\,\mathrm{bp}}}
% Supplementary numbering: Table S1, Figure S1, Section S1.
\renewcommand{\thetable}{S\arabic{table}}
\renewcommand{\thefigure}{S\arabic{figure}}
\renewcommand{\thesection}{S\arabic{section}}
% Table rows are generated by make_tables.py into tables/<name>.tex, the same
% files the manuscript uses. LaTeX's \input wrapper breaks booktabs rules inside
% a tabular; the primitive does not.
\makeatletter
\newcommand{\revrows}[1]{\@@input tables/#1.tex }
\makeatother
\setlength{\parskip}{0.5em}
\setlength{\parindent}{0pt}

\title{Supplementary Material\\[0.4em]
\large A Physics-Informed Neural Network Forward Solver for the\\
Rough-Volatility Calibration Inverse Problem}
\date{}

\begin{document}
\maketitle
\vspace{-2.5em}

This document collects the supporting tables and figures for the manuscript.
Every result here is summarised, with its key numbers, in the main text; section
references (\S) are to the manuscript. All tables are generated from the result
files in the public repository by \texttt{make\_tables.py}, and all figures by
\texttt{figures\_revision.py}.

% ------------------------------------------------------------------------------
\section{Compute environment (manuscript \S4)}
Every run reported in the manuscript and here was performed on the machine
below, with no GPU.

\begin{table}[h]
\centering
\caption{Compute environment, recorded by \texttt{env\_report.py}.}
\label{tab:S-env}
\small
\begin{tabular}{@{}l l@{}}
\toprule
Item & Value \\
\midrule
\revrows{environment}
\bottomrule
\end{tabular}
\end{table}

% ------------------------------------------------------------------------------
\section{Liquid-core refit (manuscript \S6.2)}
The full-surface fit is evaluated unchanged on the restricted points, then the
surface is refitted on those points alone from both ends of the $H$ ridge of
manuscript \S7.5. On both restricted windows the two starts converge to the same
parameters, so there $H$ is identified; on the full surface they do not.

\begin{table}[h]
\centering
\caption{Liquid-core refit. Fit quality in vol\bp, weighted by the vega weights
of the objective and unweighted. ``Full-surface fit'' rows evaluate the reference
calibration on the restricted points without refitting; the refits are labelled
by the $H$ they started from.}
\label{tab:S-liquid}
\footnotesize
\begin{tabular}{@{}l r l r r r r r@{}}
\toprule
Window & pts & Fit & weighted & unweighted & $H$ & $\rho$ & $\nu$ \\
\midrule
\revrows{liquid}
\bottomrule
\end{tabular}
\end{table}

% ------------------------------------------------------------------------------
\section{Hyperparameter sensitivity (manuscript \S7.2)}
Two seeds per configuration, $10^4$ iterations, scored against a lifted Monte
Carlo of $10^6$ antithetic paths at $1\,600$ steps. The held-out residual is the
RMS PDE residual on a fixed, seeded set of $20\,000$ collocation points never used
in training.

\begin{table}[h]
\centering
\caption{Architecture grid at $n{=}4$. Solve-error in px\bp{}; held-out PDE
residual (RMS); minutes per run on CPU. The minutes were recorded while other
jobs shared the machine and are indicative only; the controlled cost is in
manuscript Table~3.}
\label{tab:S-hpgrid}
\small
\begin{tabular}{@{}r r r r r r@{}}
\toprule
width & depth & params & solve-error & residual & min \\
\midrule
\revrows{hparams_grid}
\bottomrule
\end{tabular}
\end{table}

\begin{table}[h]
\centering
\caption{Held-out PDE residual across $n$ (width 64, depth 4,
$n_{\mathrm{col}}{=}2000$), with the solve-error and both quantities relative to
$n{=}4$. The residual is flat to within twelve per cent while the solve-error
varies by up to fifty-seven per cent.}
\label{tab:S-hpnsweep}
\small
\begin{tabular}{@{}r r r r r@{}}
\toprule
$n$ & solve-error & residual & solve/solve$_4$ & resid./resid.$_4$ \\
\midrule
\revrows{hparams_nsweep}
\bottomrule
\end{tabular}
\end{table}

\begin{table}[h]
\centering
\caption{Collocation ablation at $n{=}32$ (width 64, depth 4).}
\label{tab:S-hpcolloc}
\small
\begin{tabular}{@{}r r r r@{}}
\toprule
$n_{\mathrm{col}}$ & solve-error & residual & min \\
\midrule
\revrows{hparams_colloc}
\bottomrule
\end{tabular}
\end{table}

% ------------------------------------------------------------------------------
\section{Signed solve-error by strike (manuscript \S7.3)}

\begin{figure}[h]
\centering
\includegraphics[width=0.75\linewidth]{fig_signed.png}
\caption{Signed solve-error by strike (network minus lifted Monte Carlo, in
implied-vol basis points), seed-averaged over eight seeds, at $n{=}4,8,16,32$.
The network over-prices at and above the money at every $n$; the put-wing error
is smaller throughout and close to zero at $n{=}32$.}
\label{fig:S-signed}
\end{figure}

% ------------------------------------------------------------------------------
\section{Profile likelihood (manuscript \S7.5)}
The values plotted in manuscript Figure~4.

\begin{table}[h]
\centering
\caption{Profile likelihood in $H$ on the headline surface. $H$ is fixed at each
grid value and the other five parameters are re-optimised by Nelder--Mead from
both ends of the ridge, keeping the better. Weighted and unweighted
implied-vol RMSE in vol\bp; the last three columns show how $\nu$, $\rho$ and
$\lambda$ compensate along the ridge. From $H{=}0.30$ upward the re-optimised
$\nu$ sits on its upper bound of 2.}
\label{tab:S-profile}
\small
\begin{tabular}{@{}r r r r r r@{}}
\toprule
$H$ & weighted & unweighted & $\nu$ & $\rho$ & $\lambda$ \\
\midrule
\revrows{profile}
\bottomrule
\end{tabular}
\end{table}

% ------------------------------------------------------------------------------
\section{Risk functionals, every vector (manuscript \S7.5)}

\begin{table}[h]
\centering
\caption{Risk functionals for every parameter vector the surface cannot tell
apart, Fourier-priced at $T{=}0.25$, $S_0{=}1$, $r{=}0$. Variance swap as
$\sqrt{K_{\mathrm{var}}}$ in annualised vol points; risk reversal as
IV(25$\Delta$ call)$-$IV(25$\Delta$ put) in vol points; collar as long 90\% put,
short 110\% call, in price basis points of spot. Ridge vectors are labelled with
their degradation of the weighted objective relative to the best fit.}
\label{tab:S-risk}
\small
\begin{tabular}{@{}l l r r r@{}}
\toprule
Vector & & Var.\ swap & $25\Delta$ RR & Collar \\
\midrule
\revrows{risk}
\bottomrule
\end{tabular}
\end{table}

% ------------------------------------------------------------------------------
\section{Weekly option-implied parameters (manuscript \S7.5)}

\begin{figure}[h]
\centering
\includegraphics[width=0.92\linewidth]{fig_regime.png}
\caption{Weekly option-implied $H$ from the trade-tape panel, fitted from both
ends of the ridge (top), and which end each week prefers (bottom: loss from the
wide-box start minus loss from the canonical start; negative prefers the
smoother $H$). Shaded: January 2024, excluded as the ETF-approval window. The
week-to-week scatter and the disagreement between the two starts on the same
week are the non-identification of manuscript \S7.5 at weekly resolution; the
stability test in manuscript Table~5 compares regime means over roughly fifty
weeks each.}
\label{fig:S-regime}
\end{figure}

\end{document}
